\chapter{Compressible Flow Solver (Redesign)}
\label{chap:compressibleFlowRedesign}

\section{Synopsis}
The CompressibleFlowSolverRedesign solves the same unsteady compressible
Euler and Navier-Stokes equations in conservative variables as the
CompressibleFlowSolver described in the previous chapter -- see
Eq.~(\ref{eq:euler}) and following for the governing equations, which are
unchanged. It is built on the redesigned operator framework, in which every
step of the discretisation (basis transforms, volume fluxes, Riemann
solvers, boundary conditions, diffusion) is an operator with
interchangeable implementations for each execution space. The same session
file therefore runs

\begin{itemize}
\item in \inltt{Serial} on a single core,
\item in \inltt{AVX} using SIMD vectorisation (AVX2/AVX512 on x86, NEON on
  Arm), and
\item on a \inltt{Device} (CUDA, HIP or SYCL GPU back-ends),
\end{itemize}

and in parallel under MPI, where the trace exchange overlaps interior work
with communication. The discretisation is discontinuous Galerkin (weak DG
for the advection terms and an interior penalty scheme for the viscous
terms); a continuous projection is not available in this solver.

The solver currently provides:

\begin{itemize}
\item the Euler and Navier-Stokes equations in conservative variables
  (\inltt{EulerCFE} and \inltt{NavierStokesCFE}) with explicit time
  integration, in 1D, 2D and 3D;
\item equations of state as a templated property of every operator that
  needs one: ideal gas and Van der Waals gas;
\item Riemann solvers: \inltt{Average}, \inltt{AUSM0}--\inltt{AUSM3},
  \inltt{HLL}, \inltt{HLLC}, \inltt{LaxFriedrichs} and \inltt{Roe};
\item constant and Sutherland's law viscosity;
\item slip and no-slip (adiabatic or isothermal, including per-region wall
  temperatures) walls, farfield conditions based on Riemann invariants,
  entropy-consistent inflows, pressure outflow, stagnation inflow and
  zeroth-order extrapolation;
\item restart from checkpoint (\inltt{.fld}) files, including at a
  different polynomial order from the one the file was written at.
\end{itemize}

Features of the legacy solver that are not yet available here include the
implicit (JFNK) path, shock capturing and artificial viscosity, de-aliasing,
axisymmetric flows, quasi-1D nozzle flows and the forcing and filter
frameworks.

\section{Usage}
\begin{lstlisting}[style=BashInputStyle]
CompressibleFlowSolverRedesign session.xml
\end{lstlisting}

The execution space may be chosen on the command line without touching the
session file:

\begin{lstlisting}[style=BashInputStyle]
CompressibleFlowSolverRedesign --opExecSpace Device session.xml
\end{lstlisting}

The implementation used within a space (for example \inltt{SumFac} or
\inltt{StdMat} element kernels) can likewise be forced for every operator
at once with \inltt{--opImpl}. Both choices may equally be made from
within the session file, through the \inltt{OperatorExecSpace} and
\inltt{OperatorImplementation} properties under Solver info below, with
the command line taking precedence, which makes it convenient to run one
session file across execution spaces without editing it. When neither the
command line nor the session specifies a space, the solver defaults to
the most capable one the build supports and prints a warning saying so.

\section{Session file configuration}
\label{sec:sessionFileCompRedesign}
The session file follows the structure described in
Section~\ref{sec:sessionFileComp} for the legacy solver; below we describe
the sections that differ, and the options currently understood.

\subsection*{Parameters}
\begin{lstlisting}[style=XmlStyle]
<PARAMETERS>
  <P> TimeStep        = 0.0000001         </P>
  <P> FinTime         = 1.0               </P>
  <P> NumSteps        = FinTime/TimeStep  </P>
  <P> IO_InfoSteps    = 100               </P>
  <P> Twall           = 300.15            </P>
</PARAMETERS>
\end{lstlisting}
\begin{itemize}
\item \inltt{TimeStep} is the time-step. It is always required: the
  redesign solver integrates with a fixed step and, unlike the legacy
  solver, does not (yet) let the CFL number drive the step size.
\item \inltt{FinTime}, \inltt{NumSteps} and \inltt{IO\_InfoSteps} behave as
  in the legacy solver.
\item \inltt{Twall} sets the wall temperature for every isothermal
  (\inltt{WallViscous}) boundary at once; individual regions can override
  it as described under boundary conditions. Default value = 300.15 $K$.
\end{itemize}

\subsection*{Time integration scheme}
The \inltt{<TIMEINTEGRATIONSCHEME>} block is unchanged from the legacy
solver:
\begin{lstlisting}[style=XmlStyle]
<TIMEINTEGRATIONSCHEME>
  <METHOD> RungeKutta </METHOD>
  <ORDER> 4 </ORDER>
</TIMEINTEGRATIONSCHEME>
\end{lstlisting}
The compressible systems are integrated explicitly, so the explicit
families (\inltt{ForwardEuler}, \inltt{RungeKutta} and its SSP variants,
\inltt{AdamsBashforth}) are the relevant ones.

\subsection*{Solver info}
\begin{lstlisting}[style=XmlStyle]
<SOLVERINFO>
  <I PROPERTY="EQType"            VALUE="NavierStokesCFE"  />
  <I PROPERTY="Projection"        VALUE="DisContinuous"    />
  <I PROPERTY="AdvectionType"     VALUE="WeakDG"           />
  <I PROPERTY="DiffusionType"     VALUE="InteriorPenalty"  />
  <I PROPERTY="UpwindType"        VALUE="Roe"              />
  <I PROPERTY="ViscosityType"     VALUE="Constant"         />
</SOLVERINFO>
\end{lstlisting}
\begin{itemize}
\item \inltt{EQType} selects the equations:
  \begin{itemize}
  \item \inltt{EulerCFE} -- compressible Euler equations;
  \item \inltt{NavierStokesCFE} -- compressible Navier-Stokes equations.
  \end{itemize}
\item \inltt{Projection} must be \inltt{DisContinuous}.
\item \inltt{AdvectionType} must be \inltt{WeakDG}.
\item \inltt{DiffusionType} (Navier-Stokes only) must be
  \inltt{InteriorPenalty}.
\item \inltt{UpwindType} selects the Riemann solver for the inviscid trace
  flux: \inltt{Average}, \inltt{AUSM0}, \inltt{AUSM1}, \inltt{AUSM2},
  \inltt{AUSM3}, \inltt{HLL}, \inltt{HLLC}, \inltt{LaxFriedrichs} or
  \inltt{Roe}. This property is required.
\item \inltt{ViscosityType} (Navier-Stokes only) is \inltt{Constant}
  (default) or \inltt{Variable} for Sutherland's law.
\item \inltt{OperatorExecSpace} selects the execution space
  (\inltt{Serial}, \inltt{AVX} or \inltt{Device}) from within the session
  file; the \inltt{--opExecSpace} command-line argument overrides it.
\item \inltt{OperatorImplementation} sets the elemental implementation,
  either \inltt{SumFac} or \inltt{StdMat}, for every operator at once;
  the \inltt{--opImpl} command-line argument overrides it. When
  neither is
  given, individual operators may be assigned their own implementation
  per execution space through the \inltt{OPTIMISATION} section,
\begin{lstlisting}[style=XmlStyle]
<OPTIMISATION>
  <OPERATORS>
    <AVXBACKEND>
      <MASS> StdMat </MASS>
    </AVXBACKEND>
  </OPERATORS>
</OPTIMISATION>
\end{lstlisting}
  where \inltt{SERIALBACKEND}, \inltt{AVXBACKEND} and
  \inltt{DEVICEBACKEND} sections each list operators by name with their
  implementation. Anything still unspecified defaults to \inltt{SumFac}
  with a warning.
\item \inltt{ProjectInitialConditions}, when set to \inltt{True}, performs
  an element-local Galerkin projection of expression-based initial
  conditions instead of the default collocation interpolation, which is
  useful when comparing against a projected reference solution. The
  default (\inltt{False}) matches the legacy solver's behaviour.
\end{itemize}

\subsection*{Equation of state}
Where the legacy solver selects the equation of state through
\inltt{SOLVERINFO} and reads its constants from \inltt{PARAMETERS}, the
redesign solver has a dedicated session block, and every operator that
needs the gas law (Riemann solvers, viscous fluxes, boundary
conditions) is compiled for each equation of state:
\begin{lstlisting}[style=XmlStyle]
<EquationOfState>
  <Type> IdealGas </Type>
  <Parameters>
    <Gamma>       1.4 </Gamma>
    <GasConstant> 287 </GasConstant>
  </Parameters>
</EquationOfState>
\end{lstlisting}
Two types are available:
\begin{itemize}
\item \inltt{IdealGas} with parameters \inltt{Gamma} and
  \inltt{GasConstant};
\item \inltt{VanDerWaals} with parameters \inltt{Gamma},
  \inltt{GasConstant}, \inltt{Tcrit} and \inltt{Pcrit}.
\end{itemize}
The parameters are required -- there are no silent defaults -- and may be
expressions in terms of \inltt{PARAMETERS} entries.

\subsection*{Reference values}
The gas transport properties of the Navier-Stokes system come from a
\inltt{ReferenceValues} block rather than from the parameters:
\begin{lstlisting}[style=XmlStyle]
<ReferenceValues>
  <Viscosity>   mu     </Viscosity>
  <Prandtl>     Pr     </Prandtl>
  <Density>     rhoInf </Density>
  <Pressure>    pInf   </Pressure>
</ReferenceValues>
\end{lstlisting}
\begin{itemize}
\item \inltt{Viscosity} is the (reference) dynamic viscosity. Default
  value = 1.78e-05 $Pa\,s$.
\item \inltt{Prandtl} is the Prandtl number (default value = 0.72).
\item \inltt{thermalConductivity} may be given in its place, the Prandtl
  number then being computed from it; only one of the two may be set.
\item \inltt{Density} and \inltt{Pressure} are the reference density and
  pressure. Default values = 1.225 $Kg/m^{3}$ and 101325 $Pa$.
\item \inltt{Temperature} and \inltt{Tsutherland} are read when
  \inltt{ViscosityType} is \inltt{Variable} and set the reference and
  Sutherland temperatures of Sutherland's law. Default values = 288.15 $K$
  and 110 $K$.
\end{itemize}
Anything omitted falls back to its default with a warning, so it is good
practice to state every value the case depends on explicitly.

\subsection*{Boundary conditions}
Boundary regions and the conserved variables are declared exactly as in
the legacy solver. The following \inltt{USERDEFINEDTYPE} tags are
available; each is implemented as a boundary operator that claims the
regions carrying its tag.
\begin{center}
  \begin{tabular}{l|l}
  \toprule
  \inltt{USERDEFINEDTYPE} & Behaviour \\
  \midrule
  \inltt{Wall}, \inltt{Symmetry} & Slip (inviscid) wall / symmetry plane \\
  \inltt{WallViscous} & No-slip isothermal wall at \inltt{Twall} \\
  \inltt{WallAdiabatic} & No-slip adiabatic wall \\
  \inltt{RiemannInvariant} & Farfield in/outflow via Riemann invariants \\
  \inltt{PressureOutflow} & Outflow imposing only the static pressure \\
  \inltt{StagnationInflow} & Inflow from stagnation conditions \\
  \inltt{EnforceEntropyPressure} & Entropy-consistent inflow, pressure form \\
  \inltt{EnforceEntropyTotalEnthalpy} & Entropy-consistent inflow, total-enthalpy form \\
  \inltt{EnforceEntropyVelocity} & Entropy-consistent inflow, velocity form \\
  \inltt{ExtrapOrder0} & Zeroth-order extrapolation of the interior state \\
  \bottomrule
  \end{tabular}
\end{center}
For example, a farfield region reads:
\begin{lstlisting}[style=XmlStyle]
<REGION REF="0">
  <D VAR="rho"  USERDEFINEDTYPE="RiemannInvariant" VALUE="0" />
  <D VAR="rhou" USERDEFINEDTYPE="RiemannInvariant" VALUE="0" />
  <D VAR="rhov" USERDEFINEDTYPE="RiemannInvariant" VALUE="0" />
  <D VAR="E"    USERDEFINEDTYPE="RiemannInvariant" VALUE="0" />
</REGION>
\end{lstlisting}
A boundary tag can carry parameters of its own: an isothermal wall may be
given a temperature that differs from the session-wide \inltt{Twall} with
\begin{lstlisting}[style=XmlStyle]
<D VAR="rho" USERDEFINEDTYPE="WallViscous:twall=350.0" VALUE="0" />
\end{lstlisting}
which is how, for instance, a channel heated on one side is set up.
Ordinary Dirichlet values and periodic boundaries are also supported, the
latter across processor partitions in parallel runs.

\subsection*{Initial conditions and restarts}
Initial conditions are given either as expressions, as in the legacy
solver, or from a field file:
\begin{lstlisting}[style=XmlStyle]
<FUNCTION NAME="InitialConditions">
  <F VAR="rho,rhou,rhov,E" FILE="restart.fld" />
</FUNCTION>
\end{lstlisting}
Restart files may be read at a different polynomial order from the one
they were written at, and modal (orthogonal) as well as modified expansion
bases round-trip; the coefficients are interpolated element by element. A
warning is printed at $t=0$ if a projected initial state would be lossy at
the current order.

\section{Porting a legacy session}
A session written for the legacy \inltt{CompressibleFlowSolver} will usually
run under the redesign without complaint and give quietly wrong answers,
because several inputs are now read from different places. Each case below
leaves the session valid, the solver silent apart from a warning, and the
physics changed. Work through them in order.

\subsection*{Gas properties move to \texttt{<EquationOfState>}}
\inltt{Gamma} and \inltt{GasConstant} are no longer read from
\inltt{<PARAMETERS>}:
\begin{lstlisting}[style=XmlStyle]
<EquationOfState>
  <Type> IdealGas </Type>
  <Parameters>
    <Gamma>       Gamma       </Gamma>
    <GasConstant> GasConstant </GasConstant>
  </Parameters>
</EquationOfState>
\end{lstlisting}
The values may reference existing \inltt{<PARAMETERS>} entries by name, as
above, and leaving those parameters in place is harmless: the session reader
warns that the equation-of-state definition wins. The block itself is not
optional. Without it the operator key resolves with an empty
equation-of-state suffix and creation fails with a misleading
\begin{lstlisting}[style=BashInputStyle]
No such operator: AdvDiffTraceFluxCFERoeSerial
\end{lstlisting}
which looks like a missing Riemann solver and is not.

\subsection*{Transport properties move to \texttt{<ReferenceValues>}}
A \inltt{<ReferenceValues>} section carries everything else the legacy
solver read as a parameter; the values and their defaults are listed under
Section~\ref{sec:sessionFileCompRedesign} above.

This is the case that bites. The legacy solver read the viscosity as the
parameter \inltt{mu}; the redesign reads the reference value
\inltt{Viscosity}. A session that sets \inltt{mu = rhoInf*L*uInf/Re}
therefore gives the legacy solver whatever that evaluates to and the
redesign the built-in default -- a factor of thousands apart in a typical
case -- with nothing failing. State every value the case depends on, even
where the default happens to match, since a default that is right today is
still a silent dependency on a number defined elsewhere.

\subsection*{Read the startup warnings}
Every fallback to a default is announced, one line per missing value, among a
good deal of other output. Treat any of them as a porting error rather than as
noise:
\begin{lstlisting}[style=BashInputStyle]
CompressibleFlowSolverRedesign case.xml 2>&1 | grep "Setting default value"
\end{lstlisting}
A correctly ported viscous case prints none.

\subsection*{Schemes and forcings the redesign does not have}
\inltt{Euler}/\inltt{Forward} is not registered; use \inltt{RungeKutta} of
order 1, which is the same scheme. That one announces itself, failing with
\inltt{No such operator}.

A \inltt{<FORCING>} block does not. Forcing types are registered separately
from the solver, and a session naming one the redesign does not implement runs
to completion with the forcing simply absent, without a warning. The signature
is a partial disagreement: the quantities the forcing drives are wrong while
everything else matches. Check the block against the registered forcings, and
if the forcing is unavailable, delete it from the session rather than leave it
implying support that is not there.

\subsection*{Porting a regression test}
The interior-penalty diffusion of the redesign is not a different
discretisation from the legacy solver's, so the expectation is that a ported
case reproduces the legacy numbers to the legacy tolerances. Keep the legacy
metrics until you have proved you cannot: a gap of a few percent looks like an
acceptable discretisation difference and is more often a porting error or a
bug, and writing the redesign's own value into the \inltt{.tst} freezes that
bug in as the expected answer. If the numbers genuinely cannot be matched, say
which operator differs and why, rather than quoting a percentage.

Agreement with the legacy solver is not by itself evidence that either is
right. Where the case has an analytic handle, check it -- a wall heat flux
against $k\,\mathrm{d}T/\mathrm{d}n$ with $k = \mu c_p/\mathit{Pr}$, for
instance. That check is what exposes a viscosity read from the wrong place,
which is otherwise self-consistent: the wrong viscosity enters both sides.
Confirm too that the test can fail, by running it with the feature under test
disabled and checking that the error moves by decades.

Re-mesh when the case has outgrown its mesh, and hold two things fixed when
you do. Keep the element mix: a mesh of one element type is a weaker test,
because the boundary storage, the trace gathers and the interpolation between
differently shaped faces all vary with element type. And re-verify against the
legacy solver on the new mesh, since the reference values belong to the mesh
rather than to the case.

Finally, assert the handedness of any mesh you generate. Nektar++ does not
reject an element whose local axes are left-handed: it runs, it converges, and
it is quietly wrong. The reference prism runs $\xi_1$ along $v_0 \to v_1$,
$\xi_2$ along $v_0 \to v_3$ and $\xi_3$ along $v_0 \to v_4$, so
$(e_1 \times e_2) \cdot e_3$ must be positive; splitting a hexahedral cell
into two prisms on a diagonal and extruding, the obvious anticlockwise
winding gives exactly the wrong sign. The symptom is subtle rather than
catastrophic, so give the case a quantity that is exactly zero by symmetry to
check against.

\section{Examples}
The regression tests under
\inltt{solvers/CompressibleFlowSolverRedesign/Tests} are working examples
of every feature above, including:
\begin{itemize}
\item \inltt{IsentropicVortex.xml} -- 2D Euler, ideal gas, Lax-Friedrichs
  upwinding, periodic boundaries;
\item the \inltt{Couette\_WeakDG\_IP\_MODIFIED} cases, adiabatic and
  isothermal -- 2D Navier-Stokes with the interior
  penalty diffusion, no-slip walls, reference
  values and the equation-of-state block;
\item \inltt{RiemannInvariant\_2D\_mix.xml} and the
  \inltt{EnforceEntropy*\_2D\_mix.xml} cases -- the farfield and
  entropy-consistent boundary conditions on a mixed
  quadrilateral-triangle mesh;
\item \inltt{TGVFlow3D\_Hex8.xml} -- the 3D Taylor-Green vortex.
\end{itemize}
